\documentclass{MMCStyle}
\usepackage{MyTemplate}
\usepackage{zhlipsum}

\bianhao{MC2604019}
\tihao{ABCDEFGHI}
\timu{论文标题}
\keyword{content；content；content}
% \bibliographystyle{gbt7714-numerical}
\addbibresource{ref.bib}

\graphicspath{{./images/}}

\begin{document}
\nocite{*}
\begin{abstract}
这是摘要。

\zhlipsum[2]

摘要结束。
\end{abstract}

% 取消目录页码
\tableofcontents\newpage
\clearpage
\pagenumbering{arabic}
\pagestyle{fancy}
% 恢复默认的页面样式

$\forall a, b \in\R$ ，有 $a^2 + b^2 > 0$



\section{问题重述}


开始正文

\begin{table}[!htbp]
\centering
\caption{sjf}
\begin{tabularx}{\textwidth}{Y Y}
\toprule
阿巴阿巴 & 安居客 \\
\midrule
时间飞机 & 金石可镂 \\
sjflks & jlsk \\
\bottomrule
\end{tabularx}
\end{table}

\section{问题分析}

\zhlipsum[8]

\section{模型假设}

\zhlipsum[1]

\section{符号说明}

\zhlipsum[1]

\section{问题一模型的建立与求解}

\zhlipsum[1]

\subsection{定性分析}

\zhlipsum[1]

\subsection{解析解的求解}

\zhlipsum[10]

\subsubsection{奇异情形}

经典的 \inline{Lotka-Volterra} 模型假设捕食者与猎物共同生活在同一空间，
捕食行为使得猎物的死亡率增加和捕食者的出生率增加，得到
\[
\begin{cases}
  \displaystyle\derivative{U}{t} = r_m U - \gamma U V \\[8pt]
  \displaystyle\derivative{V}{t} = -\beta_2 V + e \gamma U V
\end{cases}
\]
式中， $r_m$ 为食草动物的 \href{https://www.zgbk.com/ecph/words?SiteID=1&ID=144328}{内禀增长率} ；
$\beta_2$ 为食肉动物的自然死亡率；
$e$ 为能量转化效率，表示每捕食一只猎物能使捕食者增加的个体数；
$\gamma$ 为捕食率。


\zhlipsum[3-4]

\section{问题二模型的建立与求解}

\zhlipsum[7]

\subsection{数学模型的建立}

\zhlipsum[8]

\begin{figure}
\centering
\includegraphics[width=0.5\textwidth]{"phase_lam0.5_mu0.4_eps2.0.pdf"}
\end{figure}

\section{问题三模型的建立与求解}

\zhlipsum[1]

% \input{<file name>}

\section{模型的优缺点分析}

\clearpage
\appendix

\section{问题一模型的建立与求解}

\zhlipsum[1]

\subsection{定性分析}

\zhlipsum[1]

\subsection{解析解的求解}

\zhlipsum[10]

\subsubsection{奇异情形}

经典的 \inline{Lotka-Volterra} 模型假设捕食者与猎物共同生活在同一空间。

\clearpage
\phantomsection
\addcontentsline{toc}{section}{参考文献}
% \bibliography{ref}
\printbibliography[title={参考文献}] % 输出列表

\end{document}
